\subsection{응용: 페이지랭크 (2026-08-31)}

페이지랭크는 행렬 계산이 정의 자체다.

\begin{itemize}
\tightlist
\item
  전이행렬 \(M\): 페이지 \(j\)가 링크 \(L_j\)개를 가지면
  \(M_{ij}=1/L_j\) (\(j\to i\) 링크가 있을 때). 열의 합이 1.
\item
  감쇠계수 \(d\) (보통 0.85)를 넣은 구글 행렬
\end{itemize}

\[
G=dM+\frac{1-d}{n}\mathbf 1\mathbf 1^{\top}
\]

\begin{itemize}
\tightlist
\item
  페이지랭크 벡터 \(\mathbf r\)은 고유값 1에 대한 고유벡터:
  \(G\mathbf r=\mathbf r\), \(\sum r_i=1\)
\item
  계산은 거듭제곱법:
  \(\mathbf r_{t+1}=dM\mathbf r_t+\frac{1-d}{n}\mathbf 1\)을 수렴할
  때까지 반복
\item
  실제 웹 규모에서는 \(M\)을 희소행렬로 저장
\end{itemize}

예: 링크 \(1\to2,\ 1\to3,\ 2\to3,\ 3\to1\), \(d=0.85\)이면
\(\mathbf r\approx(0.388,\ 0.215,\ 0.397)\). \(G\)의 나머지 고유값은
\(-0.425\pm0.425i\) (크기 약 0.60)라서 반복할 때마다 오차가 약 0.6배로
줄어든다.
